% TOP_PROBLEM: {"id": 1172, "title": "Firoozbakht's Conjecture", "category_id": 1, "category": {"id": 1, "name": "number_theory", "display_name": "Number Theory", "description": "Properties of integers, prime numbers, Diophantine equations.", "slug": "number-theory", "order_index": 1, "created_at": "2026-07-31T15:26:25.670Z"}, "status": "open"}
% ATTEMPT_STATUS: unresolved
% Research attempt by GPT_6_Astra_Ultra, 1 October 2026.
% Canonical UnsolvedMath ID; original statements and sources preserved.
% FIRST_POSTED: 2026-10-01
% Imported from: unsolvedmath_additional_500.tex; UnsolvedMath ID 1172
% !TeX program = lualatex
% Compile twice: lualatex 1172.tex
% Self-contained: no external bibliography, images, or \input files.
% Numeric section labels and visible numbers are original UnsolvedMath IDs.
\documentclass[11pt,a4paper]{article}
\usepackage[margin=24mm,headheight=14pt]{geometry}
\usepackage{fontspec}
\setmainfont{Latin Modern Roman}
\setsansfont{Latin Modern Sans}
\setmonofont{Latin Modern Mono}
\usepackage{amsmath,amssymb,mathtools}
\usepackage{unicode-math}
\setmathfont{Latin Modern Math}
\usepackage{microtype}
\usepackage{enumitem,xurl,needspace}
\usepackage[dvipsnames]{xcolor}
\usepackage{hyperref}
\hypersetup{unicode=true,colorlinks=true,linkcolor=MidnightBlue,urlcolor=MidnightBlue,
 pdftitle={UnsolvedMath Problem 1172},pdfauthor={GPT 6 Astra Ultra},bookmarksnumbered=true}
\usepackage{bookmark,tocloft,titlesec,fancyhdr}
\setlength{\cftsecnumwidth}{4.2em}
\cftsetpnumwidth{2.5em}
\cftsetrmarg{4em}
\titleformat{\section}{\Large\bfseries\raggedright}{\thesection}{0.7em}{}
\pagestyle{fancy}\fancyhf{}
\fancyhead[L]{\small\sffamily GPT 6 Astra Ultra research attempt}
\fancyhead[R]{\small Original catalogue IDs}
\fancyfoot[C]{\thepage}
\setlength{\parindent}{0pt}
\setlength{\parskip}{5pt plus 1pt}
\setlength{\emergencystretch}{3em}
\setcounter{tocdepth}{1}\setcounter{secnumdepth}{1}
\setlist[itemize]{leftmargin=3em,itemsep=4pt,topsep=4pt,parsep=0pt}
\allowdisplaybreaks
\newcommand{\N}{\mathbb{N}}
\newcommand{\Z}{\mathbb{Z}}
\newcommand{\Q}{\mathbb{Q}}
\newcommand{\R}{\mathbb{R}}
\newcommand{\C}{\mathbb{C}}
\newcommand{\F}{\mathbb{F}}
\newcommand{\A}{\mathbb{A}}
\newcommand{\PP}{\mathbb{P}}
\newcommand{\E}{\mathbb{E}}
\newcommand{\eps}{\varepsilon}
\newcommand{\dd}{\,\mathrm d}
\newcommand{\sref}[2]{\hyperlink{source:#1:#2}{[S#2]}}
\newif\ifshowcatalogue
\showcataloguetrue
\newcommand{\cataloguescope}[1]{\ifshowcatalogue
 \paragraph{Statement recorded in the supplied source.}
 #1\fi}
\begin{document}
% UnsolvedMath ID 1172; source catalogue code NT-027

\setcounter{section}{1171}

\section{Firoozbakht's Conjecture}\label{1172}

{\small\sffamily UnsolvedMath ID 1172 · Number Theory}

\subsection{Definitions and mathematical statement}

\paragraph{Shared notation.}Unless a section says otherwise, $\N=\{1,2,\ldots\}$,
$\Z$ is the ring of integers, $\Q,\R,\C$ are the rational, real, and complex
fields, $[N]=\{1,\ldots,\lfloor N\rfloor\}$, and $\log$ is the natural
logarithm. For real functions of a parameter tending to infinity,
$f\ll g$ and $f=O(g)$ mean $|f|\leq Cg$ eventually for some fixed $C$;
$f\gg g$ reverses this comparison; $f=o(g)$ means $f/g\to0$;
$f\sim g$ means $f/g\to1$; and $f\asymp g$ means both order bounds.
A subscript on $O$ or $\ll$ specifies allowed constant dependence.
The expression $n^{o(1)}$ in an upper bound means growth smaller than
$n^\epsilon$ for each fixed $\epsilon>0$. Local definitions govern any
exception, finite-field convention, or use of the same symbol for another
object. An asymptotic claim is not automatically an effective algorithm.



Let $p_n$ be the $n$th prime. The proposed strict decrease is $p_{n+1}^{1/(n+1)}<p_n^{1/n}$ for every $n\geq1$. Equivalently, $p_{n+1}^n<p_n^{n+1}$, which avoids any choice of logarithm or numerical approximation. \sref{1172}{1} \sref{1172}{2}

\paragraph{Statement recorded in the supplied source.}
Is the sequence $p_n^{1/n}$ strictly decreasing, where $p_n$ is the $n$-th prime? \sref{1172}{1}

\subsection{Short English statement}

Does taking the nth root of the nth prime always produce a strictly decreasing sequence?

\subsection{Research attempt}

\paragraph{Scope and route.}
The inequality is strict and is asserted for every consecutive-prime
index, including $n=1$. Visser's primary paper [S2], version 2, reports
a large finite verification; its abstract was inspected on 1 October
2026. That published computation is not rerun or extended here. The
present route translates the root inequality into an exact gap threshold,
derives elementary sufficient and necessary inequalities, and checks
which information an asymptotic estimate would actually supply.

\paragraph{The exact gap threshold.}
Set $p=p_n$, $g=p_{n+1}-p_n$, and $L=\log p$. Taking logarithms,
the desired assertion is equivalent successively to
\[
 n\log(p+g)<(n+1)\log p,
 \qquad n\log(1+g/p)<L,
\]
and hence to
\[
 g<T_n:=p\bigl(e^{L/n}-1\bigr).
\]
All these equivalences preserve strictness. Since $p_{n+1}$ and $p_n$
are different primes, equality $p_{n+1}^n=p_n^{n+1}$ is also impossible
by unique factorization. Integer exponentiation is therefore a clean
way to test the original comparison without rounding logarithms.

\paragraph{An explicit sandwich requiring no prime theorem.}
For $t>0$, integration of $1/(1+x)$ over $[0,t]$ gives
\[
 \frac{t}{1+t}<\log(1+t)<t.
\]
The upper bound proves the sufficient condition
\[
                     g\le\frac{pL}{n}.
\]
The weak sign is legitimate here because the logarithmic inequality
is strict for positive $g$. Conversely the conjectured inequality,
together with the lower bound, implies
\[
 g<\frac{pL}{n-L}\qquad\text{whenever }n>L.
\]
Thus the target lies between two explicitly specified gap conditions.
Neither replaces the actual prime gap by an average gap. For large
$n$ their separation is approximately $pL^2/n^2$, but a comparison
with the actual integer gap still needs rigorous bounds in the relevant
direction.

For a uniform sufficient criterion involving only $p$, suppose one
has an explicit upper bound $n=\pi(p)\le B(p)$. Then
$pL/n\ge pL/B(p)$, and $g\le pL/B(p)$ suffices. An upper bound on
$\pi(p)$ is the useful direction for this particular sufficient
condition; substituting a lower bound would make the permitted gap
larger and would not justify the conclusion. Any use of a published
prime-counting estimate must track its threshold and its direction.

\paragraph{The lower-order scale is visible in a conditional expansion.}
To separate the algebra from the input, assume an estimate of the form
\[
 \pi(p)=\frac pL\left(1+\frac1L+\frac2{L^2}
                    +O(L^{-3})\right).
\]
Then inversion of the power series gives
\[
 \frac{pL}{n}=L^2-L-1+O(L^{-1}).
\]
Also $e^x-1=x+O(x^2)$ for $x=L/n\to0$, so
\[
 T_n=\frac{pL}{n}+O\!\left(\frac{pL^2}{n^2}\right)
     =L^2-L-1+O(L^{-1}).
\]
The last simplification uses $n\asymp p/L$, making the exponential
remainder $O(L^4/p)$. This is a conditional algebraic consequence of
the stated estimate, not a gap theorem. Even an eventual upper bound
$g\le L^2$ would not, by itself, imply $g<T_n$, because it misses
the negative term of size $L$. Likewise $g=O(L^2)$ leaves its implied
constant and the needed lower-order saving unspecified.

\paragraph{Why a prime-counting asymptotic cannot force this locally.}
One can construct increasing real sequences with the same first-order
growth as the primes but with isolated jumps violating this sort of
threshold. Let $N_j=2^{2^j}$ and, for integers $n\ge2$, put
\[
 a_n=n\log n+\sum_{N_j\le n}\sqrt{N_j}.
\]
Both terms are nondecreasing and the first strictly increases. Because
the $N_j$ grow doubly exponentially, the accumulated jump sum is
$O(\sqrt n\log\log(n+2))=o(n\log n)$; hence $a_n\sim n\log n$.
At an index with $n+1=N_j$, however,
$a_{n+1}-a_n\ge\sqrt{N_j}$. The threshold analogous to $T_n$ for
$a_n$ is $O((\log n)^2)$, as follows by substituting
$a_n\asymp n\log n$ into its definition. For sufficiently large $j$
the jump exceeds that threshold.

These are not primes and do not provide a counterexample. They show
that first-order growth information alone cannot prove a comparison of
two adjacent roots. A proof must use additional arithmetic control of
the actual consecutive prime gaps, sufficiently strong at every index.

\paragraph{Executed strict comparisons.}
The script generates consecutive primes by an integer sieve and checks
\[
 p_{n+1}^{\,n}<p_n^{\,n+1}\qquad(1\le n\le1000)
\]
with arbitrary-precision integer powers. All pass; the final tested
pair is $(7919,7927)$. At $n=1$, the comparison is $3<4$, so the
initial case is included explicitly. No floating-point approximation
or comparison of rounded logarithms is used. See
\texttt{checks.py} and \texttt{results.json} in
\path{_Local/Erdos_problems/UnsolvedMath/attempt_audit_2026_10_01/number_theory/}.
This is a small independent check, not the large computation in [S2].

\paragraph{Remaining gap and outcome.}
The exact gap threshold and the one-sided sufficient and necessary
bounds are proved. What is missing is a uniform estimate for the actual
gap below that threshold, including its lower-order terms where needed,
or a single violating prime pair. The asymptotic calculations supply
neither. \textbf{Outcome: UNRESOLVED.} In particular no deduction from
an average prime spacing to a pointwise bound is made.

\subsection{Sources}

{\small

\begin{itemize}

\item[\hypertarget{source:1172:1}{S1}] ULAM.AI, \emph{UnsolvedMath}, supplied \texttt{problems.json}, record 1172 (NT-027), statement and background. The embedded literature-review date is 2026-08-17; that is the archive’s review date, not a fresh certification. Dataset landing page: \url{https://huggingface.co/datasets/ulamai/UnsolvedMath}.

\item[\hypertarget{source:1172:2}{S2}] Verifying the Firoozbakht, Nicholson, and Farhadian conjectures up to the 81st maximal prime gap, arXiv:1904.00499. \url{https://arxiv.org/abs/1904.00499}. Bibliographic information imported from the supplied record.

\item[\hypertarget{source:1172:3}{S3}] Firoozbakht's conjecture overview. \url{https://en.wikipedia.org/wiki/Firoozbakht%27s_conjecture}. Bibliographic information imported from the supplied record.


\item[E1] Matt Visser, \emph{Verifying the Firoozbakht,
Nicholson, and Farhadian conjectures up to the 81st maximal prime gap},
arXiv:1904.00499v2 (2019).
\url{https://arxiv.org/abs/1904.00499v2}. Abstract inspected
1 October 2026; its stated finite verification is background only.

\end{itemize}

}

\label{end:1172}
\end{document}
